\chapter{External Tool Comparison: Reproducibility Protocol}
\label{app:external-comparison}

\paragraph{Scope.}
This appendix collects the full reproduction protocol for the head-to-head
comparison between PROTEA and eggNOG-mapper~v2.1.13 reported in
\cref{subsec:exp7-eggnogmapper}, targeting reviewers who wish to verify,
extend, or replicate the evaluation independently. It covers test-set
extraction, tool installation, run parameters, evaluation alignment with
the CAFA protocol, and complete per-cell $F_{\max}$ tables that were
abbreviated in the main text. It does not overlap with
\cref{app:api} (REST API reference) or \cref{app:adr-log} (architecture
decision records): those appendices address system design, whereas this one
addresses empirical reproducibility of a specific comparative experiment.

This appendix documents the exact procedure used to evaluate eggNOG-mapper
on the same test set and evaluation protocol as PROTEA (\cref{subsec:exp7-eggnogmapper}).
It is intended to allow full reproduction of the comparison results and to
serve as a template for evaluating additional external tools.

% ────────────────────────────────────────────────────────────
\section{Test Set Extraction}
\label{app:ext-test-set}

The test set consists of the 20{,}281 proteins that gained at least one new
experimental GO annotation between GOA version~220 and version~229 (the same
temporal holdout used for all PROTEA experiments in \cref{chap:evaluation}).
The FASTA file was extracted from PROTEA's database using the built-in API
endpoint:

\begin{verbatim}
curl -o test_proteins.fasta \
  "http://localhost:8000/annotations/evaluation-sets/\
42b34e79-6fe9-4fa0-b718-02f43a1e3192/delta-proteins.fasta"
\end{verbatim}

Each FASTA header contains the UniProt accession, entry name, organism,
taxonomy ID, and CAFA category (NK/LK/PK):

\begin{verbatim}
>A0A024RBG1 NUD4B_HUMAN OS=Homo sapiens OX=9606 (LK)
MMKFKPNQTRTYDREGFKKRAACLCFRSEQEDEVLLVSSS...
\end{verbatim}

The resulting file contains 20{,}281 sequences spanning all three CAFA categories
(NK: 2{,}831, LK: 3{,}410, PK: 15{,}313).

% ────────────────────────────────────────────────────────────
\section{eggNOG-mapper Setup}
\label{app:ext-emapper-setup}

eggNOG-mapper v2.1.13 was run via its official Biocontainers Docker image
to ensure a reproducible environment. The following steps were performed:

\paragraph{1.\ Pull the Docker image.}
\begin{verbatim}
docker pull quay.io/biocontainers/\
  eggnog-mapper:2.1.13--pyhdfd78af_2
\end{verbatim}

\paragraph{2.\ Download the eggNOG databases (v5.0.2).}
At the time of this experiment (March 2026), the official download script
pointed to a deprecated URL. The databases were downloaded manually from
\texttt{eggnog5.embl.de}:

\begin{verbatim}
# eggnog.db (6.3 GB compressed, ~39 GB uncompressed)
wget -O eggnog.db.gz \
  http://eggnog5.embl.de/download/emapperdb-5.0.2/\
eggnog.db.gz
gunzip eggnog.db.gz

# Diamond protein database (4.8 GB compressed, ~8.7 GB)
wget -O eggnog_proteins.dmnd.gz \
  http://eggnog5.embl.de/download/emapperdb-5.0.2/\
eggnog_proteins.dmnd.gz
gunzip eggnog_proteins.dmnd.gz

# Taxonomy database (~71 MB compressed)
wget -O eggnog.taxa.tar.gz \
  http://eggnog5.embl.de/download/emapperdb-5.0.2/\
eggnog.taxa.tar.gz
tar xzf eggnog.taxa.tar.gz
\end{verbatim}

Total disk usage: approximately 48~GB.

\paragraph{3.\ Verify installation.}
\begin{verbatim}
docker run --rm \
  -v /path/to/eggnog-data:/data \
  quay.io/biocontainers/\
    eggnog-mapper:2.1.13--pyhdfd78af_2 \
  emapper.py --version --data_dir /data

# Output:
# emapper-2.1.13 / eggNOG DB version: 5.0.2
# Diamond version: 2.0.15 / MMseqs2: 16.747c6
\end{verbatim}

% ────────────────────────────────────────────────────────────
\section{Running eggNOG-mapper}
\label{app:ext-emapper-run}

\begin{verbatim}
docker run --rm \
  -v /path/to/eggnog-data:/data \
  -v /path/to/input:/input \
  -v /path/to/output:/output \
  quay.io/biocontainers/\
    eggnog-mapper:2.1.13--pyhdfd78af_2 \
  emapper.py \
    -i /input/test_proteins.fasta \
    --data_dir /data \
    -o test_proteins \
    --output_dir /output \
    --cpu 8 \
    -m diamond \
    --go_evidence experimental \
    --tax_scope auto \
    --target_orthologs all
\end{verbatim}

\paragraph{Key parameters.}
\begin{description}
  \item[\texttt{-m diamond}] Use DIAMOND for sequence search (default mode).
  \item[\texttt{--go\_evidence experimental}] Transfer only GO terms supported
    by experimental evidence codes, matching the evidence code filter used in
    PROTEA's evaluation protocol.
  \item[\texttt{--tax\_scope auto}] Automatically determine the taxonomic scope
    for ortholog assignment.
  \item[\texttt{--target\_orthologs all}] Consider all orthologs (one-to-one,
    one-to-many, many-to-many).
  \item[\texttt{--cpu 8}] Use 8 CPU threads for the DIAMOND search phase.
\end{description}

\paragraph{Runtime.} The full run completed in approximately 21 minutes
on the same workstation used for all PROTEA experiments (64~GB RAM, no GPU
required for eggNOG-mapper).

\paragraph{Output.} eggNOG-mapper produced annotations for 19{,}958 of the
20{,}281 input proteins (98.4\% hit rate). Of those, 17{,}334 (85.5\% of the
test set) received at least one GO term prediction. The remaining 2{,}947
proteins either had no DIAMOND hit or no GO terms were transferred from the
matched orthologous group.

% ────────────────────────────────────────────────────────────
\section{CAFA Evaluation Protocol}
\label{app:ext-cafa-eval}

To ensure a fair comparison, eggNOG-mapper predictions were evaluated using
exactly the same protocol as all PROTEA experiments:

\begin{enumerate}
  \item \textbf{Ground truth computation.} The NK/LK/PK classification and
    ground-truth annotations were computed from the same EvaluationSet
    (GOA~220$\to$229) using PROTEA's \texttt{compute\_evaluation\_data()}
    function, which implements NOT-qualifier propagation through the GO DAG
    following the CAFA5 protocol.
  \item \textbf{Prediction format.} eggNOG-mapper does not produce per-term
    confidence scores; each predicted GO term is assigned a uniform score of
    1.0 (binary: predicted or not predicted).
  \item \textbf{Evaluation.} The \texttt{cafaeval} tool was run with the same
    parameters as all other experiments: GO propagation mode \texttt{max},
    normalisation mode \texttt{cafa}, and IA weighting from
    \texttt{IA\_cafa6.tsv}.
\end{enumerate}

The evaluation script \texttt{scripts/evaluate\_external\_tool.py} automates
this entire pipeline. It accepts the external tool's output file, connects
to PROTEA's database to retrieve the ground truth, and runs \texttt{cafaeval}
for all three CAFA settings (NK, LK, PK):

\begin{verbatim}
poetry run python scripts/evaluate_external_tool.py \
  --evaluation-set-id \
    42b34e79-6fe9-4fa0-b718-02f43a1e3192 \
  --tool emapper \
  --input output/test_proteins.emapper.annotations
\end{verbatim}

% ────────────────────────────────────────────────────────────
\section{Raw Results}
\label{app:ext-raw-results}

\Cref{tab:emapper-raw} presents the complete per-cell results for
eggNOG-mapper alongside the PROTEA methods.

\begin{table}[htbp]
  \centering
  \caption[Complete Fmax comparison]{Complete $F_{\max}$ comparison: eggNOG-mapper vs.\ PROTEA methods (IA-weighted).}
  \label{tab:emapper-raw}
  \resizebox{\textwidth}{!}{%
  \begin{tabular}{l ccc ccc ccc}
    \toprule
    & \multicolumn{3}{c}{\textbf{NK}} & \multicolumn{3}{c}{\textbf{LK}} & \multicolumn{3}{c}{\textbf{PK}} \\
    \cmidrule(lr){2-4} \cmidrule(lr){5-7} \cmidrule(lr){8-10}
    \textbf{Method} & BPO & MFO & CCO & BPO & MFO & CCO & BPO & MFO & CCO \\
    \midrule
    eggNOG-mapper 2.1.13 & 0.247 & 0.359 & 0.386 & 0.382 & 0.334 & 0.450 & 0.190 & 0.199 & 0.325 \\
    PROTEA baseline      & 0.412 & 0.590 & 0.668 & 0.467 & 0.558 & 0.675 & 0.187 & 0.278 & 0.325 \\
    PROTEA full-feature  & \textbf{0.431} & \textbf{0.620} & \textbf{0.692} & \textbf{0.478} & \textbf{0.607} & \textbf{0.697} & \textbf{0.201} & \textbf{0.297} & \textbf{0.339} \\
    \bottomrule
  \end{tabular}
  }
\end{table}

\Cref{tab:emapper-delta} shows the absolute $F_{\max}$ difference between
PROTEA's full-feature re-ranker and eggNOG-mapper.

\begin{table}[htbp]
  \centering
  \caption[Absolute Fmax improvement vs.\ eggNOG-mapper]{Absolute $F_{\max}$ improvement: PROTEA full-feature vs.\ eggNOG-mapper.}
  \label{tab:emapper-delta}
  \begin{tabular}{lccc}
    \toprule
    \textbf{Category} & \textbf{BPO} & \textbf{MFO} & \textbf{CCO} \\
    \midrule
    NK & +0.184 & +0.261 & +0.306 \\
    LK & +0.096 & +0.273 & +0.247 \\
    PK & +0.011 & +0.098 & +0.014 \\
    \bottomrule
  \end{tabular}
\end{table}

\paragraph{Coverage analysis.} eggNOG-mapper found DIAMOND hits for 19{,}958
proteins but only transferred GO terms to 17{,}334 (85.5\% of the test set).
PROTEA, by contrast, produces predictions for all 20{,}281 proteins (100\%
coverage) because every protein has an embedding and therefore has $k$
nearest neighbours. This coverage gap accounts for part of the $F_{\max}$
difference, particularly in NK where 14.5\% of proteins receive no
eggNOG-mapper prediction.

% ────────────────────────────────────────────────────────────
\section{Methodological Notes}
\label{app:ext-notes}

\paragraph{Score assignment.}
eggNOG-mapper produces a binary set of GO terms per protein without
per-term confidence scores. For the CAFA evaluation, all predicted terms are
assigned a confidence of 1.0. This means the precision-recall curve is a
single point (one threshold), and $F_{\max}$ equals the $F_1$ score at that
point. Methods that produce graded confidence scores (such as PROTEA) can
optimise the precision-recall trade-off by varying the threshold, which may
contribute to higher $F_{\max}$ values.

\paragraph{GO propagation.}
eggNOG-mapper transfers GO terms from orthologous groups but does not
explicitly propagate them up the GO DAG. The \texttt{cafaeval} evaluator
applies \texttt{prop=max} propagation during evaluation, so ancestor terms
are included in the comparison for all methods.

\paragraph{Evidence code filter.}
The \texttt{--go\_evidence experimental} flag restricts eggNOG-mapper to
transferring GO terms supported by experimental evidence codes in the eggNOG
database. This is conceptually aligned with PROTEA's use of experimental-only
annotations from GOA, though the exact set of supporting evidence and the
database versions differ.
